\documentclass[dvipsnames]{article}

\usepackage[utf8]{inputenc}
\usepackage{fontspec}
\usepackage{polyglossia}
\usepackage{unicode-math}
\usepackage{csquotes}
\setdefaultlanguage{russian}
\setotherlanguage{english}

\setmainfont{CMU Serif}
\setsansfont{CMU Sans Serif}
\setmonofont{CMU Typewriter Text}
% \setmathfont{Latin Modern Math}

\usepackage[letterpaper,top=2cm,bottom=2cm,left=3cm,right=3cm,marginparwidth=1.75cm]{geometry}

\usepackage{amsmath}
\usepackage{graphicx}
\usepackage{amsthm}
\usepackage{array}
\usepackage{fancyhdr}
\usepackage{bussproofs}
\usepackage[colorlinks=true, allcolors=blue]{hyperref}
\usepackage[backend=biber]{biblatex}
\usepackage{xcolor}

\usepackage{minted2}
\usemintedstyle{xcode}
\setminted{escapeinside=||}
\newminted[stella]{cpp}{escapeinside=||,frame=single}
\newmintinline[haskell]{haskell}{fontsize=\scriptsize}
\newmintinline[java]{java}{fontsize=\scriptsize}
\newmintinline[python]{python}{fontsize=\scriptsize}
\newmintinline[code]{c}{fontsize=\small}
\newmintinline[smallcode]{c}{fontsize=\small}
\let\smallcodeOrig\smallcode
\renewcommand{\smallcode}[2][]{%
  \ifmmode
    \if\relax\detokenize{#1}\relax
      \mathchoice{\mbox{\normalfont\small\ttfamily #2}}{\mbox{\normalfont\small\ttfamily #2}}{\mbox{\normalfont\small\ttfamily #2}}{\mbox{\normalfont\small\ttfamily #2}}%
    \else
      \mathchoice{\mbox{\normalfont\small\ttfamily #2}}{\mbox{\normalfont\small\ttfamily #2}}{\mbox{\normalfont\small\ttfamily #2}}{\mbox{\normalfont\small\ttfamily #2}}%
    \fi
  \else
    \smallcodeOrig[#1]{#2}%
  \fi}
\let\codeOrig\code
\renewcommand{\code}[2][]{%
  \ifmmode
    \if\relax\detokenize{#1}\relax
      \mathchoice{\mbox{\normalfont\small\ttfamily #2}}{\mbox{\normalfont\small\ttfamily #2}}{\mbox{\normalfont\small\ttfamily #2}}{\mbox{\normalfont\small\ttfamily #2}}%
    \else
      \mathchoice{\mbox{\normalfont\small\ttfamily #2}}{\mbox{\normalfont\small\ttfamily #2}}{\mbox{\normalfont\small\ttfamily #2}}{\mbox{\normalfont\small\ttfamily #2}}%
    \fi
  \else
    \codeOrig[#1]{#2}%
  \fi}

\usepackage{sourcecodepro}
\usepackage{mdframed}

\DeclareMathOperator{\xor}{\mathsf{xor}}
\DeclareMathOperator{\tru}{\mathsf{tru}}
\DeclareMathOperator{\fls}{\mathsf{fls}}

\addbibresource{../references.bib}

\begin{document}

\pagestyle{fancy}
\fancyhf{}
\chead{Задание выполнил \textcolor{blue}{Бенжамин Пирс (замените на своё имя)}}
\cfoot{Языки программирования. Семантика и системы типов (весна 2026, Университет ИТМО)}

\section*{Тема 7. Теоретическое задание \textcolor{blue}{(выполнил Бенжамин Пирс)}}%

\begin{enumerate}
  \item Рассмотрите \code{error} и \code{try}-\code{with} как \emph{производную форму}, реализующую обработку ошибок через типы-суммы~\cite[\S~11.9]{TaPL-ru}.
    \begin{itemize}
      \item \code{error} ведёт себя как \code{inl unit};
      \item константы (например, \code{0}) оборачиваются в \code{inr} (например, \code{inr 0});
      \item все прочие конструкции явно обрабатывают тип \code{Unit + T} вместо выражения типа \code{T}; например, для применения $\textcolor{red}{e}(\code{t$_1$ t$_2$})$:
        \begin{minted}[escapeinside=??,frame=single]{text}
case t?$_1$? of
  inl u => inl u
| inr t?$_1'$? =>
    case t?$_2$? of
      inl u => inl u
    | inr t?$_2'$? => inr (?$\textcolor{red}{e}($?t?$_1')$? ?$\textcolor{red}{e}($?t?$_2')$?)
        \end{minted}
    \end{itemize}
    \begin{enumerate}
      \item Выпишите явно функцию раскрытия сокращений для $\lambda$-исчисления
        с логическими и арифметическими выражениями~\cite[\S\S~8--9]{TaPL-ru}
        и со \emph{всеми} простыми расширениями из~\cite[\S~11]{TaPL-ru}.
        Если для какого-то расширения \textbf{невозможно} раскрытие, сохраняющее семантику исходных правил вычисления
        для \code{error} и \code{try}-\code{with}~\cite[\S~14.1]{TaPL-ru}, то кратко объясните, почему.
      \item Покажите, что раскрытие сохраняет семантику исходных правил вычисления для \code{error} и \code{try}-\code{with}~\cite[\S~14.1]{TaPL-ru}.
      Для доказательства достаточно рассмотреть только простое типизированное $\lambda$-исчисление с логическими~\cite[\S\S~8--9]{TaPL-ru}, без арифметических выражений и дополнительных расширений.
    \end{enumerate}

  \color{blue}
  \textbf{Решение.}
  \emph{место для решения}
  \color{black}

  \clearpage

  \item Рассмотрите расширение простого типизированного $\lambda$-исчисления с обработкой ошибок со значениями~\cite[\S~14.3]{TaPL-ru} конструкцией
    \[
      \code{with exception}~\mathcal{E}~\code{do}~t_1~\code{catch}~x \Rightarrow t_2.
    \]
    Эта конструкция:
    \begin{itemize}
      \item вводит тип исключений $\code{T}_\mathsf{exn} = \mathcal{E}$ \emph{временно}, только на время вычисления $t_1$;
      \item при возникновении ошибки (\code{raise}) управление передаётся ветке
        $\code{catch}~x \Rightarrow t_2$, где $x$~--- переменная, с которой связывается значение ошибки типа $\mathcal{E}$;
      \item в $t_2$ допускается использование переменной $x$.
    \end{itemize}
    \begin{enumerate}
      \item Выпишите явно правила вычисления для этой конструкции.
      \item Предложите правила типизации для этой конструкции, а также новые (уточнённые) правила типизации \code{raise} и \code{try}-\code{catch},
        принимающие во внимание локальный тип исключений. При необходимости расширьте отношение типизации
        (например, добавьте текущий тип исключений в контекст или введите суждение вида $\Gamma \vdash_{\mathcal{E}} t : T$).
      \item Аргументируйте, почему предложенные правила сохраняют безопасность типизации,
        или укажите причину, по которой предложенная конструкция не может быть безопасной для системы типов.
        Подробное доказательство безопасности типизации не требуется.
    \end{enumerate}

  \color{blue}
  \textbf{Решение.}
  \emph{место для решения}
  \color{black}
\end{enumerate}

\section*{Список литературы}
\printbibliography[heading=none]

\end{document}
